% GENERATED FILE. Edit canonical lessons, then run npm run generate:books.
% canonical-source-hash: fnv1a64:ac064000eb317b85
% canonical-lessons: IT-C446-mura, IT-C446-rocca, IT-C446-borgo, IT-C446-chiostro, IT-C446-vagare

\chapter{Il centro storico}
\label{ch:it-il-centro-storico}
\hlchaptermodality{446}

\begin{chapteropening}
\textbf{By the end of this chapter:} \emph{I can describe an old town: the walls, the fortress, the village, the cloister, and wandering round.}

Complete the last lesson of chapter 446: I can describe an old town: the walls, the fortress, the village, the cloister, and wandering round.
\end{chapteropening}

\section[le mura]{le mura — the walls (of a town)}

\label{lesson:IT-C446-mura}



\begin{quote}
Before the new one: say the Italian for the arcade, then the Italian for the loggia.
\end{quote}

\subsection*{You'll want to know: le mura}
\textbf{le mura} — \textquotedblleft{}the walls (of a town)\textquotedblright{}.

\begin{quote}
Possiamo camminare sulle mura della città.
\end{quote}

\begin{quote}
\emph{We can walk on the city walls.}
\end{quote}

\emph{Il muro} is a wall; the old plural \emph{le mura} means the walls round a town.

\begin{cousinweb}
\emph{mura} is the old plural of \emph{muro} 'wall', which is inherited from Latin \emph{mūrus} (EtymDB).
\end{cousinweb}

\begin{grammarlens}[title={What you've built}]
One more word for telling, explaining and describing at B1.
\end{grammarlens}

\subsection*{Your turn}
\begin{itemize}
\raggedright
  \item \emph{Say it:} \emph{le mura}
  \item \emph{Say it:} \emph{le mura}, once more
  \item \emph{Say it:} a sentence of your own with \emph{le mura}
  \item \emph{Recall:} say the Italian for the arcade, then the Italian for the loggia, then say \emph{le mura} again
\end{itemize}

\begin{tcolorbox}[breakable,colback=teal!4,colframe=teal!35!black,title={Before you move on}]
What does le mura mean? (\textquotedblleft{}The walls (of a town)\textquotedblright{}.) Say it once more.
\end{tcolorbox}
\section[la rocca]{la rocca — the fortress, the stronghold}

\label{lesson:IT-C446-rocca}



\begin{quote}
Before the new one: say the Italian for the loggia, then the Italian for the walls.
\end{quote}

\subsection*{You'll want to know: la rocca}
\textbf{la rocca} — \textquotedblleft{}the fortress, the stronghold\textquotedblright{}.

\begin{quote}
La rocca domina il paese.
\end{quote}

\begin{quote}
\emph{The fortress looks down on the village.}
\end{quote}

Usually on a rock above the town.

\begin{cousinweb}
\emph{rocca} goes back to Vulgar Latin \emph{rocca} 'rock, crag'; the deeper origin is not settled, and it may be older than Latin.
\end{cousinweb}

\begin{grammarlens}[title={What you've built}]
One more word for telling, explaining and describing at B1.
\end{grammarlens}

\subsection*{Your turn}
\begin{itemize}
\raggedright
  \item \emph{Say it:} \emph{la rocca}
  \item \emph{Say it:} \emph{la rocca}, once more
  \item \emph{Say it:} a sentence of your own with \emph{la rocca}
  \item \emph{Recall:} say the Italian for the loggia, then the Italian for the walls, then say \emph{la rocca} again
\end{itemize}

\begin{tcolorbox}[breakable,colback=teal!4,colframe=teal!35!black,title={Before you move on}]
What does la rocca mean? (\textquotedblleft{}The fortress, the stronghold\textquotedblright{}.) Say it once more.
\end{tcolorbox}
\section[il borgo]{il borgo — the village, the old quarter}

\label{lesson:IT-C446-borgo}



\begin{quote}
Before the new one: say the Italian for the walls, then the Italian for the fortress.
\end{quote}

\subsection*{You'll want to know: il borgo}
\textbf{il borgo} — \textquotedblleft{}the village, the old quarter\textquotedblright{}.

\begin{quote}
Il borgo ha solo trecento abitanti.
\end{quote}

\begin{quote}
\emph{The village has only three hundred inhabitants.}
\end{quote}

A small old town, often with walls.

\begin{cousinweb}
\emph{borgo} comes from Late Latin \emph{burgus} 'small fort' (EtymDB), which was borrowed from a Germanic word for a fort.
\end{cousinweb}

\begin{grammarlens}[title={What you've built}]
One more word for telling, explaining and describing at B1.
\end{grammarlens}

\subsection*{Your turn}
\begin{itemize}
\raggedright
  \item \emph{Say it:} \emph{il borgo}
  \item \emph{Say it:} \emph{il borgo}, once more
  \item \emph{Say it:} a sentence of your own with \emph{il borgo}
  \item \emph{Recall:} say the Italian for the walls, then the Italian for the fortress, then say \emph{il borgo} again
\end{itemize}

\begin{tcolorbox}[breakable,colback=teal!4,colframe=teal!35!black,title={Before you move on}]
What does il borgo mean? (\textquotedblleft{}The village, the old quarter\textquotedblright{}.) Say it once more.
\end{tcolorbox}
\section[il chiostro]{il chiostro — the cloister}

\label{lesson:IT-C446-chiostro}



\begin{quote}
Before the new one: say the Italian for the fortress, then the Italian for the village.
\end{quote}

\subsection*{You'll want to know: il chiostro}
\textbf{il chiostro} — \textquotedblleft{}the cloister\textquotedblright{}.

\begin{quote}
Nel chiostro c'è un pozzo antico.
\end{quote}

\begin{quote}
\emph{In the cloister there is an old well.}
\end{quote}

A square garden with arcades, inside a monastery.

\begin{cousinweb}
\emph{chiostro} is inherited from Latin \emph{claustrum} 'bar, enclosure' (EtymDB), from \emph{claudere} 'to shut'.
\end{cousinweb}

\begin{grammarlens}[title={What you've built}]
One more word for telling, explaining and describing at B1.
\end{grammarlens}

\subsection*{Your turn}
\begin{itemize}
\raggedright
  \item \emph{Say it:} \emph{il chiostro}
  \item \emph{Say it:} \emph{il chiostro}, once more
  \item \emph{Say it:} a sentence of your own with \emph{il chiostro}
  \item \emph{Recall:} say the Italian for the fortress, then the Italian for the village, then say \emph{il chiostro} again
\end{itemize}

\begin{tcolorbox}[breakable,colback=teal!4,colframe=teal!35!black,title={Before you move on}]
What does il chiostro mean? (\textquotedblleft{}The cloister\textquotedblright{}.) Say it once more.
\end{tcolorbox}
\section[vagare]{vagare — to wander, to roam}

\label{lesson:IT-C446-vagare}



\begin{quote}
Before the new one: say the Italian for the village, then the Italian for the cloister.
\end{quote}

\subsection*{You'll want to know: vagare}
\textbf{vagare} — \textquotedblleft{}to wander, to roam\textquotedblright{}.

\begin{quote}
Abbiamo vagato per il borgo tutto il pomeriggio.
\end{quote}

\begin{quote}
\emph{We wandered round the village all afternoon.}
\end{quote}

With no fixed goal; \emph{vago} — vague.

\begin{cousinweb}
\emph{vagare} is inherited from Latin \emph{vagārī} 'to wander, to stroll about' (EtymDB), from \emph{vagus} 'roaming'.
\end{cousinweb}

\begin{grammarlens}[title={What you've built}]
One more word for telling, explaining and describing at B1.
\end{grammarlens}

\subsection*{Your turn}
\begin{itemize}
\raggedright
  \item \emph{Say it:} \emph{vagare}
  \item \emph{Say it:} \emph{vagare}, once more
  \item \emph{Say it:} a sentence of your own with \emph{vagare}
  \item \emph{Recall:} say the Italian for the village, then the Italian for the cloister, then say \emph{vagare} again
\end{itemize}

\begin{tcolorbox}[breakable,colback=teal!4,colframe=teal!35!black,title={Before you move on}]
What does vagare mean? (\textquotedblleft{}To wander, to roam\textquotedblright{}.) Say it once more.
\end{tcolorbox}
